\documentclass[11pt]{article}
\usepackage[T1]{fontenc}
\usepackage{lmodern,microtype}
\usepackage{amsmath,amssymb,amsthm}
\usepackage[letterpaper,margin=1in]{geometry}
\usepackage{booktabs,array,enumitem}
\usepackage[colorlinks=true,linkcolor=blue!50!black,urlcolor=blue!50!black]{hyperref}
\newcommand{\GL}{\mathrm{GL}}
\newcommand{\A}{\mathbb A}
\newcommand{\tH}{\widetilde H}
\newcommand{\Ad}{\operatorname{Ad}}
\theoremstyle{plain}
\newtheorem*{lem*}{Lemma}
\newtheorem*{obs*}{Observation}
\title{Referee report on \texttt{mjl\_part4\_draft5.tex}, with suggestions
for the wild integrality problem}
\author{Claude (Anthropic), for David Joyner}
\date{November 7, 2026}
\begin{document}
\maketitle

\section*{0. Scope}
\paragraph{Read in full.} Draft~5 (1589 lines; 22 pages; it compiles
cleanly).

\paragraph{Compiled numbering of the new results.}
\begin{itemize}
\item Example~6.3 (data);
\item Lemma~6.4 (toral Fourier test);
\item Proposition~6.5 (all quadratic tori);
\item \S6.4 (Steps~T and~G);
\item Proposition~8.3 (rank-two spherical unit);
\item Remark~8.4;
\item Proposition~8.5 (conditional comparison with $\Lambda$);
\item Remark~8.7 (identity germs);
\item Example~8.8 (numerical ledger).
\end{itemize}

\paragraph{Versions and sources.}
\begin{itemize}
\item Draft~5 cites Parts~I and~II as Draft~8 and Part~III as Draft~9. I
have Draft~6 of each, and of [JC] I have revised12 but did not read it.
\item I could not verify the following against a source:
\begin{itemize}
\item [JC], Proposition~2.4 (Stone--von Neumann);
\item the exact pointwise form of [FK, \S12, Corollary];
\item Li's real-place hypothesis.
\end{itemize}
\item The literature I suggest in \S3 is cited from memory and should be
checked.
\end{itemize}

\section*{Summary verdict}
\begin{itemize}
\item \emph{\S\S2--5, 7.} Unchanged from Draft~3 apart from
the points I raised, all of which are handled. One exception: the
$u(x)$ paragraph now says ``$u(x)=1$ by definition of the relevant
convention''. That is acceptable if [FK, \S4] defines it so, but it is
no longer an argument. Keep the one-line argument from Draft~3 as well.
\item \emph{The wild rank-two example (\S\S6.3--6.4).} This is a genuine
and correct new computation. I rechecked every value and every
decomposition (\S1.1). There is one misstatement in the description of
$B_d$.
\item \emph{Proposition~8.5.} It contains a new hypothesis, that
nongood ordinary classes have no contribution. This hypothesis is
unnecessary, and it rests on a mismatched pairing of classes (\S1.2).
It should be deleted. As written it also makes the proposition hard to
apply with unit elements at almost all places.
\item \emph{Proposition~8.3, Remark~8.7, Example~8.8.} Correct. One
citation is unverified (\S1.3).
\end{itemize}

\paragraph{The main recommendation.} Add the positivity lemma of \S2.
It is elementary, and it strengthens Proposition~6.1 considerably:
\emph{for the unitary matrix $(a_{\rho\tau})$, integrality is
equivalent to nonnegativity}, i.e.\ to every $L_D\Theta_\rho$ being a
positive-definite function. This changes what has to be proved, and it
makes several of the routes in \S3 much shorter.

\section*{1. Correctness of the new material}

\subsection*{1.1 The rank-two wild example: correct, one slip}
I verified the following.
\begin{itemize}
\item \emph{The setup.} $F=\mathbb Q_3(\zeta_3)$, $|3|_F=1/9$. $D=E_0\oplus E_0j$ with $j^2=\varpi$ is
division, since $\varpi$ has odd valuation and is not a norm from the
unramified $E_0=F(i)$. $E_0$ is unramified because $-1$ is not a square
in $\mathbb F_3$.
\item \emph{The unramified torus, $R_{E_0}=|a^2+b^2|/|3a^2-b^2|$.} Both
values are right. $a^2+b^2$ is a unit because $-1$ is a nonsquare. For
$b=\varpi u$, $3a^2-b^2=3(a^2+\zeta_3u^2)$ with $a^2+u^2\not\equiv0$.
So $R\in\{9,1\}$ as stated, and
$R_{E_0}=1+2\,\mathrm{reg}_{C_4}=3+2\chi_1+2\chi_2+2\chi_3$.
\item \emph{The ramified tori.} $R_{E_d}=|a^2-db^2|/|3a^2+db^2|$. The
three valuation cases give $9,3,1$. The decomposition
$1+\mathrm{reg}_{C_2}+\mathrm{reg}_{C_6}$ has the values $9$, $3$ and
$1$ on the identity, the even nonidentity elements and the odd elements.
The Fourier coefficients $(3,2,1,1,1,1)$ are right.
\item \emph{Lemma~6.4.} Correct.
\begin{itemize}
\item $\zeta_{\alpha\chi}$ is well defined, using injectivity of $P_n$ on
$A_E$ and $\theta(b)=\tilde\theta(b^*)$.
\item It is multiplicative, using multiplicativity of the star on the
torus.
\item Stone--von Neumann characters are $b_D\zeta$ on $C_E$ and $0$
off it.
\item The restriction identity follows.
\end{itemize}
\emph{Repair:} Lemma~6.4 has no status label; it should be \textbf{[P]}.
\item \emph{Proposition~6.5.} Correct, including the choice of
$\tilde\theta$ with $\tilde\theta(a^*)=1$ (the star image meets $\mu_n$
trivially) and the Galois invariance.
\end{itemize}

\paragraph{The slip.} In the description of
$B_d=E_d^\times/F^\times U_1$, the text says that even classes ``modulo
$F^\times$ are represented by $1+cj_d$ with $c\in\varpi\mathcal O_F$''.
Those elements lie in $U_1$, so they are trivial in $B_d$. It should say
$c\in\mathcal O_F$, read modulo $\varpi$. The multiplication law and the
additivity of $\bar c$ are then right. It would also help to note that
$U_1=\mathcal O_F^\times(1+\mathfrak p_{E_d}^2)$.

\paragraph{Consistency check.} The torus results are exactly
positive-definiteness of the restriction of $L_D\mathbf 1$ to each
covered torus. This is consistent with \S2 below, and it is a necessary
condition.

\subsection*{1.2 Proposition~8.5: delete the nongood-class hypothesis}
The proposition assumes that every \emph{nongood} rational noncentral
ordinary class contributes zero, ``removing the ordinary classes without
a covered counterpart''. But in the comparison the ordinary class $[s]$
is paired with the covered class of $s^*=\iota(s)^n$, not with
$\iota(s)$.
\begin{itemize}
\item This is what hypotheses (i)--(ii) say: the local identities
compare $J(s_v^*,\tilde f_v)$ with $J(s_v,f_v)$.
\item It is also what the proof itself uses, via $\iota(s^n)=\iota(s)^n$.
\item The star of any element is good, since
$[\tilde t^n,\tilde u]=[\tilde t,\tilde u]^n=1$.
\item Theorem~5.2 gives a bijection from \emph{all} noncentral
$[s]\in E^\times/K^\times$ onto the good classes.
\end{itemize}
So every ordinary class has a counterpart, and no extra support
condition is needed.

As written, the hypothesis would require the ordinary test function to
kill every $[s]$ with $s\notin K^\times E^{\times n}$. That is
incompatible with unit elements at almost all places, unless the
function is adjusted at some place to vanish on specific classes.
\emph{Delete the hypothesis} and the corresponding sentence of the proof
and ledger. ``Nongood ordinary class'' is meaningful only for the
lift $\iota(\gamma)$, which is not the partner of $\gamma$.

\subsection*{1.3 Other new material}
\paragraph{Proposition~8.3.} The elementary part is correct:
\begin{itemize}
\item $J(\mathrm{diag}(a,b),1_{K_0})=1$ for $a,b$ units;
\item both sides vanish when the valuations differ;
\item scalar translates contribute the central character.
\end{itemize}
The final sentence relies on [FK, \S12, Corollary] giving a
\emph{pointwise} identity $J(x^*,\tilde f^0)=J(x,f^0)$. I could not
verify that form.

For $r\ge3$, Lemma~8.2 shows that such a pointwise identity would force
$J(\zeta x,1_{K_0})=J(x,1_{K_0})$ on mixed tori. That is a checkable
consequence. Either it confirms the pointwise form for $r=3$, or it
shows that FK's identity involves a root sum. Remark~8.4 should mention
this test rather than assert that FK ``cover'' mixed tori.

\paragraph{Remark~8.7 and Example~8.8.} Correct and well hedged. The
ledger arithmetic is right: $b=3$, $\Lambda=1$ forces
$\lambda_\infty=1/3$, and $(1/3)/9=1/27$ balances $27$.

\paragraph{Ledger.} $N_2=N_3=1$ and $N_5=4$ are right.

\section*{2. A positivity lemma (recommended for the paper)}
\begin{lem*}[\textbf{[P]}; strengthens Proposition~6.1]
Let $(a_{\rho\tau})$ be the matrix of $L_D$ in the orthonormal character
bases (Part~III, Theorem~4.3 and Corollary~5.1). The following are
equivalent:
\begin{enumerate}[label=(\alph*)]
\item $L_D$ permutes the irreducible characters;
\item every $a_{\rho\tau}$ is an integer;
\item every $a_{\rho\tau}$ is real and $\ge0$;
\item every $L_D\Theta_\rho$ is a positive-definite function on $\tH$.
\end{enumerate}
Moreover, for a single $\rho$, condition (c) for that row (equivalently
(d) for that $\rho$) already implies that $L_D\Theta_\rho$ is a
nonnegative combination of characters with norm~$1$.
\end{lem*}

\begin{proof}
\emph{(c)$\Rightarrow$(a).} Let $\rho\ne\rho'$. Orthogonality of rows
gives $\sum_\tau a_{\rho\tau}a_{\rho'\tau}=0$, a sum of nonnegative
terms. So distinct rows have disjoint supports. Since $L_D$ is unitary,
each column has norm~$1$. A column meets the support of only one row,
so its unique nonzero entry is~$1$. Each row has norm~$1$ and its
nonzero entries are~$1$, so it has exactly one of them.

\emph{(a)$\Rightarrow$(b)$\Rightarrow$(c).} (a)$\Rightarrow$(b) is
trivial. (b)$\Rightarrow$(c) is Proposition~6.1.

\emph{(c)$\Leftrightarrow$(d).} On a group compact modulo the center,
with a fixed unitary central character, a continuous class function is
positive definite if and only if its coefficients in the irreducible
characters are~$\ge0$ (by Peter--Weyl and GNS).
\end{proof}

\paragraph{Consequences.}
\begin{enumerate}[label=(\roman*)]
\item \emph{Integrality is a positivity problem.} Any construction
exhibiting $L_D\Theta_\rho$ as the character of an actual, possibly
reducible, representation proves irreducibility.
\item \emph{A cheaper finite test.} At a finite level, convolution by
$L_D\Theta_\rho$ acts on the $\tau$-isotypic part of the twisted group
algebra by the scalar $a_{\rho\tau}/\dim\tau$, up to the normalization
of measures. So it suffices to check that the Hermitian matrix
$\bigl(L_D\Theta_\rho(h^{-1}g)\bigr)_{g,h}$ is positive semidefinite.
No character table of the cover is needed.
\item \emph{A global argument needs only positivity.} In the global
argument of \S3.5, the exact value of $\Lambda$ becomes irrelevant; only
its positivity matters.
\end{enumerate}

\section*{3. Suggestions for the wild integrality problem in general}
These are research suggestions, ordered roughly by how directly they
address the problem. I mark what I have checked and what is heuristic.

\subsection*{3.1 A clean formula for $R_E$ and a counting identity (checked)}
Write $u_{ij}=\lambda_i/\lambda_j$, where the $\lambda_i$ are the
eigenvalues of $t$. For a field torus $|u_{ij}|=1$, so
$\Delta(t)=\prod_{i<j}|1-u_{ij}|$, with no square roots. Hence
\[
R_E(t)=\prod_{i<j}\frac{|1-u_{ij}|}{|1-u_{ij}^n|}
=\prod_{i<j}\ \prod_{\zeta\in\mu_n\smallsetminus1}|1-\zeta u_{ij}|^{-1},
\]
using $1-X^n=\prod_{\zeta}(1-\zeta X)$, which holds because
$\mu_n\subset F$.

\paragraph{The case $r=2$.} Put $u=t/\sigma(t)$. By Hilbert~90 this
identifies $A_E$ with $E^1$, and
\[
R_E(u)=\prod_{\zeta\ne1}|1-\zeta u|^{-1}.
\]
If $\zeta$ has order prime to $p$, then $|1-\zeta u|=1$. Indeed,
$u\equiv\zeta^{-1}$ would give $N(u)\equiv\zeta^{-2}$, so $\zeta^2=1$.
Hence only the $p$-power roots of unity in $\mu_n$ contribute.

So for $r=2$, Step~T reduces to comparing $|1-u|$ with $|1-\zeta|$ for
$\zeta\in\mu_{p^a}$, along the unit filtration of $E^1$. The standard
behavior of $x\mapsto x^p$ on $U^i$ changes at $i=e/(p-1)$. This should
allow Step~T to be proved for \emph{all} wild $F$ and all $n$ when
$r=2$, with $R_E$ written explicitly as a combination of indicators of
the subgroups $E^1\cap U_E^k$. I recommend doing this next, since it
turns Proposition~6.5 from an example into a theorem.

\paragraph{A counting identity (any $r$).} Let $\Lambda\subset\mathfrak
d/\mathfrak e$ be an $\Ad(t)$-stable lattice, and $N\gg0$. Then
\[
\Delta(t)^{-2}=\#\mathrm{Fix}\bigl(\Ad t\ \text{on}\ \Lambda/\varpi^N\Lambda\bigr),
\qquad
R_E(t)^2=\frac{\#\mathrm{Fix}(\Ad t^n)}{\#\mathrm{Fix}(\Ad t)}.
\]
Proof: $|\det(1-\Ad t)|^{-1}=[\Lambda:(1-\Ad t)\Lambda]$, and kernel and
cokernel have equal order on the finite module. So $R_E^2$ is a ratio
of permutation-character values. This is the natural form in which to
compare $R_E$ with the Frobenius formula of \S3.3.

\subsection*{3.2 Gluing via Brauer's characterization (heuristic, but standard in shape)}
At a finite level, the question is whether a class function on a finite
twisted group is a virtual character. Brauer's characterization theorem
says this is detected by restrictions to \emph{elementary} subgroups:
cyclic $p'$-groups times $p$-groups. By Proposition~6.1, once
$L_D\Theta_\rho$ is a virtual character it is irreducible, because it
has norm~$1$ and is positive at~$1$.

So Step~G reduces to elementary subgroups. The $p'$-parts are handled
by tame (Adams-type) arguments. The core is the restriction to the
covered pro-$p$ part, $\widetilde{U_D^1}$ together with the $p$-part of
the covered torus. The twisted version of Brauer's theorem should be
located in the literature (e.g.\ Karpilovsky's books on projective
representations) or proved. I have not checked that it exists in the
form needed.

\subsection*{3.3 Explicit construction: tame types for $D^\times$ and their covers}
Since $(n,r)=1$, a wild place for $n$ has $p\ne r$. So $D$ is
\emph{tamely} ramified, and every irreducible representation of
$D^\times$ comes from Howe's admissible pairs $(E,\xi)$, with $E/F$ of
degree~$1$ or~$r$. Explicit character formulas are given in Corwin--Howe
(Pacific J.\ Math.\ 73, 1977), and degrees in Corwin--Moy--Sally.

\paragraph{Proposal.}
\begin{enumerate}[label=(\arabic*)]
\item \emph{Define genuine admissible pairs} $(E,\tilde\xi)$, with
$\tilde\xi(zt^*)=\tilde\theta(z)\xi(t)$ on the center $C_E$ of
$\widetilde{E^\times}$.
\item \emph{Build $\tau(E,\tilde\xi)$} as an induced representation
$\mathrm{Ind}_{\tilde J}^{\tH}\tilde\Lambda$, as in Howe's construction.
Use the Stone--von Neumann representation of $\widetilde{E^\times}$ on
the torus part, and the compact splitting $\kappa^*$ on deep congruence
subgroups, over which the cover splits.
\item \emph{Compute $\Theta_\tau$ on elliptic elements by the Frobenius
formula} and compare with $L_D\Theta_{\rho(E,\xi)}$. The fixed-point
form of $R_E^2$ in \S3.1 is what such a comparison will produce.
\end{enumerate}
This is the route of Bushnell--Henniart's explicit essentially tame
Jacquet--Langlands correspondence, and of Zou's simple types for covers
(which assume $p\nmid n$). The division-algebra case is the simplest
one: no Levi subgroups, and only simple strata. By \S2 it suffices to
show that each $L_D\Theta_\rho$ is the character of \emph{some}
representation.

\subsection*{3.4 Kirillov theory for the deep part (heuristic, with a dimension check)}
On a deep congruence subgroup the cover splits. There, Howe's Kirillov
theory for compact $p$-adic groups (Pacific J.\ Math.\ 73, 1977), or the
orbit method for saturable pro-$p$ groups (Gonz\'alez-S\'anchez), writes
characters as Fourier transforms of coadjoint orbits. In exponential
coordinates the $n$th-power map is the dilation $X\mapsto nX$. So $L_D$
should act by \emph{dilating orbits}. Dilation sends orbits to orbits,
which would give irreducible characters directly, and $R_E$ would be
the Jacobian correction.

\paragraph{Dimension check.} A generic orbit of $D^\times$ has
dimension $r(r-1)$. Dilating by $n^{-1}$ multiplies its volume by
$|n|^{-r(r-1)}$, so it multiplies the degree by $|n|^{-r(r-1)/2}$. The
Heisenberg factor on the torus contributes $b_D$. Together:
\[
b_D\,|n|_F^{-r(r-1)/2}=|n|_F^{-(r^2-1)/2}=c_D,
\]
exactly the predicted degree ratio. This makes the orbit method a
promising mechanism for the pro-$p$ part in Step~G. The applicability
conditions (saturability, $p$ relative to the dimension) need checking.

\subsection*{3.5 The global route (conditional; needs only positivity)}
\begin{enumerate}[label=(\arabic*)]
\item Choose the global $D$ with $S$ containing \emph{every} place above
$n$, together with the target wild place $w$. Then every finite split
place is tame.
\item Take unit elements at the split finite places. This needs [FK,
\S12] under $(n,N_r)=1$, which is automatic for $r\le5$ with $n$ odd,
plus transport to the cover of Part~I.
\item At the tame places of $S$, Part~III, Theorem~7.6, gives
permutation matrices.
\item Globalize $\rho$ (as in [FK, \S20]). Take
$\tilde f_w=e_\tau$ (a projector), so that $f_w=T^{-1}e_\tau$ has
$\mathrm{tr}\,\rho(f_w)=a_{\rho\tau}$.
\item Separate by the Hecke algebra at split places, and use strong
multiplicity one for $D^\times$ (Badulescu). This should give
\[
a_{\rho\tau}=(\text{positive constant})\times(\text{nonnegative integer}).
\]
\end{enumerate}
By \S2, this already proves integrality, \emph{whatever the constant is}.
The remaining inputs are those of Proposition~8.5:
\begin{itemize}
\item existence of the split-place and complex-place matching, with
positive constants;
\item the germ issue at regular $x$ with $x^n$ singular.
\end{itemize}
Note also: with $S\supseteq\{v\mid n\}$, the product formula gives
$\prod_{v\in S}b_{D,v}=\prod_{v\mid\infty}|n|_v^{(r-1)/2}$. So
$\Lambda=1$ is equivalent to $\lambda_v=|n|_v^{-(r-1)/2}$ at each
complex place (compare Example~8.8). But by \S2 this exact value is not
needed for integrality.

\subsection*{3.6 Order of work}
\begin{enumerate}[label=(\alph*)]
\item Add the positivity lemma (\S2), and fix \S1.1 and \S1.2.
\item Run the finite PSD test of \S2(ii) for
$(2,3,\mathbb Q_3(\zeta_3))$, $\rho=\mathbf1$. This is cheaper than a
character table.
\item Prove Step~T for all $r=2$, using the formula of \S3.1.
\item Attack Step~G for $r=2$ by \S3.3, the explicit types. This case
is small: $E$ is quadratic, and the strata are short.
\end{enumerate}

\section*{4. Editorial}
\begin{itemize}
\item Lemma~6.4 needs a status label.
\item The abstract is now long, about 280 words; trim it.
\item The acknowledgments sentence ``added in Draft~5 after a separate
comparison \ldots'' is drafting history. Move it to the log.
\item Restore the $\{\pm1\}\cap\mu_n$ argument alongside ``by
convention'' (Summary).
\end{itemize}

\end{document}
